\chapter{The Law of Evidence Through Bayes}
\label{ch:evidence}

\section{Rules about information}

The law of evidence is a set of rules about information: which items a fact finder may hear, how they must be presented, and what directions the fact finder must be given about them. Many of these rules were developed for juries, and many have been relaxed as trials by judge alone became more common. They remain the core of every evidence course, and each can be read as a response to a problem of inference. Some rules exclude evidence whose likelihood ratio is close to one but which fact finders tend to overweight. Some require corroboration of evidence with a high or correlated error rate. Some govern the reliability of expert evidence, which supplies likelihood ratios that fact finders cannot assess for themselves. This chapter works through the main doctrines in English and American law and asks what a Bayesian account says about each.\index{evidence, law of}

The organising idea is simple. Let $E$ be an item of evidence and $H$ the proposition in issue. Evidence affects the fact finder's posterior through its likelihood ratio, $\LR = \Prob(E\given H)/\Prob(E\given\bar H)$. A fact finder may err in two ways: by misjudging the likelihood ratio, for instance treating a ratio of two as if it were twenty, or by misjudging the prior. Exclusionary rules, warnings and directions are tools for reducing these errors. A rule is justified, in decision terms, if the expected loss from errors with the rule is lower than without it, counting the information lost by exclusion.

\section{Relevance and probative value}

Chapter~\ref{ch:odds} showed that the legal definitions of relevance can be read in Bayesian terms: evidence is relevant if its likelihood ratio differs from one. Rule 403 of the Federal Rules of Evidence permits a court to exclude relevant evidence if its probative value is substantially outweighed by a danger of unfair prejudice, confusing the issues, misleading the jury, undue delay, wasting time or needlessly presenting cumulative evidence. English law has a similar general discretion in civil cases under the Civil Procedure Rules and, in criminal cases, section 78 of the Police and Criminal Evidence Act 1984 allows exclusion of prosecution evidence whose admission would have such an adverse effect on the fairness of the proceedings that it ought not to be admitted.\index{Rule 403}\index{probative value}

Probative value has a natural measure: the expected weight of evidence, or more usefully, the expected reduction in decision loss from hearing the evidence. Prejudice also has a natural measure: the expected increase in decision loss from the fact finder's misjudgment of the likelihood ratio. Suppose the true likelihood ratio is $\lambda$ but the jury treats it as $\lambda^\kappa$ with $\kappa > 1$, overweighting it. On the log scale the distortion is $(\kappa-1)\log\lambda$, which is proportional to the true weight of the evidence and vanishes as $\lambda\to 1$. In this model the distortion matters when $\kappa$ is large and the prior is close enough to the threshold that the true posterior falls below it while the distorted posterior rises above it. A Bayesian Rule 403 balance compares the value of information of the evidence, computed with the true $\lambda$, with the expected loss from distortion, computed with the jury's $\lambda^\kappa$. If instead the overweighting is an additive shift on the log scale, as in treating a ratio of two as if it were twenty, the distortion does not vanish as $\lambda\to 1$, and evidence with a small true likelihood ratio and a large tendency to be overweighted should be excluded.

\section{Character and bad character}

\subsection{The common law rule}

At common law the prosecution could not lead evidence of the defendant's bad character, such as previous convictions, to show that he was the kind of person likely to commit the offence charged. The rule had exceptions, notably for similar fact evidence whose probative force derived from striking similarity. In England, the Criminal Justice Act 2003 replaced the common law with a statutory scheme: bad character evidence is admissible through one of seven gateways, including where it is relevant to an important matter in issue between the defendant and the prosecution, such as his propensity to commit offences of the kind charged. Where it is admitted under the gateways for an important matter in issue or an attack on another person's character, the court must, on the defendant's application, exclude it if its admission would have such an adverse effect on the fairness of the proceedings that it ought not to be admitted. In \emph{R v Hanson} the Court of Appeal gave guidance on propensity evidence, emphasising that a single previous conviction will often not show propensity, and that the court should consider the similarity and number of previous offences and the time since \citep{hanson2005}.\index{bad character}\index{R v Hanson}

\subsection{The Bayesian analysis}

Is a previous conviction for burglary relevant to a charge of burglary? Let $C$ be the event that the defendant has a previous burglary conviction. The likelihood ratio is $\Prob(C\given G)/\Prob(C\given\bar G)$, where $G$ is guilt of the present burglary. The numerator is the proportion of burglars who have previous burglary convictions; the denominator is the proportion of people in the pool of possible offenders, who are innocent of this burglary, with such convictions. Because burglars reoffend at high rates and the general population has few burglary convictions, the ratio can be large, perhaps ten or more. Propensity evidence is often highly probative.

The case for excluding it, or admitting it only with caution, rests on three further considerations. First, the relevant comparison is not with the general population but with the pool of suspects, and if the police find suspects by looking at people with previous convictions, the denominator is the proportion of those suspects with convictions, which is high by construction. The likelihood ratio of a conviction among suspects identified by their record is close to one, because the record is then the reason the person is a suspect and adds little beyond it. This is a statement about how the pool was selected, not a dilution of the evidence by the search, which Chapter~\ref{ch:bf} shows does not occur for a database match. Second, juries may overweight character evidence, treating propensity as near proof. Third, admitting such evidence may encourage police to round up the usual suspects rather than investigate. The \emph{Hanson} guidance, which asks about the number and similarity of previous offences, is in effect asking whether the likelihood ratio is large among the relevant pool of suspects.

\section{Hearsay}

\subsection{The rule and its exceptions}

Hearsay is an out of court statement offered to prove the truth of what it asserts. At common law it was inadmissible unless within an exception, because the maker of the statement could not be cross examined, and so the fact finder could not test the four testimonial capacities: perception, memory, narration and sincerity. In England the Criminal Justice Act 2003 relaxed the rule in criminal proceedings, admitting hearsay in many circumstances subject to safeguards, and the Civil Evidence Act 1995 effectively abolished it in civil proceedings. In the United States the Federal Rules of Evidence retain the rule with many exceptions, and the Confrontation Clause of the Sixth Amendment limits the use of testimonial hearsay against criminal defendants, as the Supreme Court held in \emph{Crawford v Washington} \citep{crawford2004}.\index{hearsay}

\subsection{A Bayesian account}

A statement is a signal about the fact it asserts, with a likelihood ratio that depends on the declarant's perception, memory, narration and sincerity. Cross examination is a way of learning about these parameters, reducing the fact finder's uncertainty about the likelihood ratio. Without cross examination, the fact finder must integrate over its uncertainty about the declarant's reliability. Schum's analysis of the credibility of testimony decomposes the credibility of a report into three attributes, veracity, objectivity and observational sensitivity, which enter its likelihood ratio and overlap with these capacities \citep{schum1994}. Hearsay is not worthless, but its likelihood ratio is uncertain, and the uncertainty is wider when the declarant is unknown, the circumstances unrecorded and the motive unexamined.

The hearsay exceptions can be read as categories in which the likelihood ratio is likely to be large and reliably estimated without cross examination. Statements against interest are unlikely to be false, since people rarely make false statements that harm themselves. Excited utterances are made before there is time to fabricate. Business records are made routinely by people with no motive to falsify. Dying declarations, the oldest exception, were thought reliable because a dying person would not meet his maker with a lie, an assumption about sincerity that a modern court might doubt. Each exception is a prior about the reliability of a class of statements.

\section{Identification evidence}

\subsection{The Turnbull guidelines}

Mistaken eyewitness identification is a leading cause of wrongful convictions. In \emph{R v Turnbull} the Court of Appeal in 1976 laid down guidelines for cases depending wholly or substantially on disputed identification \citep{turnbull1977}. The judge must warn the jury of the special need for caution, explain that a mistaken witness can be a convincing one, and direct the jury to examine closely the circumstances of the identification: how long the witness observed the accused, at what distance, in what light, whether the view was impeded, whether the witness had seen the accused before, how long elapsed between the observation and the identification, and whether there were material discrepancies between the description and the accused. Where the quality of the identification is poor, the judge should withdraw the case unless there is supporting evidence.\index{Turnbull guidelines}\index{identification evidence}

\subsection{Identification as a noisy signal}

An eyewitness identification is a signal whose likelihood ratio depends on the conditions of observation. Let $I$ be the event that the witness identifies the defendant in a line up. Under guilt, $\Prob(I\given G)$ is the hit rate, which varies with conditions. Under innocence, with a fair line up of $k$ members, $\Prob(I\given\bar G)$ is the probability that the witness picks the innocent defendant, which is the false identification rate divided by $k$ if the line up is fair, and larger if the defendant stands out. Laboratory and field studies have estimated hit and false identification rates under different conditions. A good identification under good conditions, with a fair line up, may have a likelihood ratio of ten or more. A fleeting glimpse at night, with a suggestive procedure, may have a likelihood ratio close to one. The \emph{Turnbull} factors are the covariates that determine the likelihood ratio, and the guidelines are a verbal regression model.

The guidelines also illustrate the problem of confidence. Witnesses who are confident are believed more, but confidence at trial is weakly related to accuracy, because it grows with repetition and feedback \citep{wells1998}. Confidence at the time of the first identification, under uncontaminated procedures, is more informative \citep{wixted2017}. A Bayesian treatment would weight the identification by the confidence expressed at the first, uncontaminated procedure, not by the confidence displayed in court.

\section{Confessions}

Confessions are powerful evidence because people rarely confess to crimes they did not commit. Rarely is not never. The Innocence Project reports that between a quarter and nearly a third of the wrongful convictions in the United States later overturned by DNA evidence involved false confessions, and \citet{garrett2010} analyses the 40 false confessions among the first 250 DNA exonerations, many obtained after long interrogations of vulnerable suspects. The Police and Criminal Evidence Act 1984 excludes confessions obtained by oppression or in consequence of anything said or done likely to render them unreliable. The Kercher case of Chapter~\ref{ch:stories} involved statements obtained in a night interrogation that the European Court of Human Rights later found had violated the suspect's rights.\index{confessions}

The Bayesian analysis of a confession is that its likelihood ratio is $\Prob(\text{confess}\given G)/\Prob(\text{confess}\given\bar G)$, and the denominator depends on the interrogation. Long, isolating and coercive interrogations raise the probability that innocent people confess, sometimes dramatically, while raising the probability that the guilty confess less. A confession obtained under such conditions has a much smaller likelihood ratio than one volunteered. A confession that contains details only the perpetrator could know has a large likelihood ratio, provided the details were not supplied by the interrogators, which is why the recording of interrogations matters: it lets the court check whether the details originated with the suspect.

\section{Corroboration}

Historically English law required, or required a warning about the absence of, corroboration for certain kinds of evidence: the evidence of accomplices, of children, and of complainants in sexual cases. The requirements for children's evidence were removed by the Criminal Justice Act 1988, and the mandatory warnings for accomplices and complainants were abolished by the Criminal Justice and Public Order Act 1994, leaving a discretion. Scots law still requires corroboration in criminal cases as a general rule, so that no one may be convicted on the evidence of a single source, a requirement that has been the subject of repeated proposals for reform.\index{corroboration}

A corroboration requirement says that the posterior must be supported by at least two independent items. In Bayesian terms, it is a crude safeguard against overweighting a single item whose likelihood ratio is uncertain. If the fact finder's error in assessing a single item's likelihood ratio is large, requiring two independent items reduces the chance that a single overweighted item drives the decision. The cost is that some true cases supported by one strong item cannot be proved. The value of the requirement depends on how badly fact finders misjudge single items and how often true cases rest on one. Independence matters: two items that share a source, such as two accounts derived from one informant, are not corroboration in the Bayesian sense, as the Italian courts recognised for cooperating witnesses (Chapter~\ref{ch:gould}).

\section{Expert evidence}

\subsection{Admissibility}

In the United States, \emph{Daubert v Merrell Dow Pharmaceuticals} made the trial judge a gatekeeper for scientific evidence, directing courts to consider whether a theory or technique can be and has been tested, whether it has been subjected to peer review, its known or potential error rate, the existence of standards controlling its operation, and its general acceptance \citep{daubert1993}. Rule 702 was amended in 2000 to codify \emph{Daubert}, and again in 2023 to emphasise that the proponent must show it is more likely than not that the expert's opinion reflects a reliable application of reliable principles and methods. In England and Wales, the Criminal Procedure Rules and associated practice directions set out factors for assessing the reliability of expert evidence, following a Law Commission report.\index{Daubert}\index{expert evidence}

\subsection{Error rates}

Of the \emph{Daubert} factors, the known or potential error rate is the most directly Bayesian. An expert's conclusion is a signal, and its likelihood ratio is determined by its error rates. A forensic method that reports a match with probability $0.99$ when the samples share a source and $0.01$ when they do not supplies a likelihood ratio of ninety nine for a reported match. If the method's error rates are unknown, its likelihood ratio is highly uncertain, and a report of a match may be worth much less than it appears. The report of the President's Council of Advisors on Science and Technology in 2016 emphasised that many forensic feature comparison methods, such as bite mark analysis and firearms analysis and footwear comparisons, lacked empirical studies establishing their error rates, and argued that such studies were a precondition of scientific validity \citep{pcast2016}. In Bayesian terms, without validated error rates the likelihood ratio is highly uncertain and, as with any uncertain likelihood ratio, the fact finder should integrate over that uncertainty and discount the report heavily.

\subsection{Expert opinion on the ultimate issue}

Courts have long been wary of experts who express opinions on the ultimate issue, such as whether the defendant is guilty. The Bayesian framework explains why. The expert is qualified to assess the likelihood ratio of the evidence within her expertise, but not the prior or the other evidence, which belong to the fact finder. An expert who states a posterior probability of guilt has implicitly assumed a prior, usurping the fact finder's role. The reporting standards for forensic science that recommend likelihood ratios rather than posterior statements, discussed in Chapter~\ref{ch:odds}, and Fienberg and Kadane's principles for presenting Bayesian analyses, discussed in Chapter~\ref{ch:bf}, implement this division.

\section{Privileges and exclusionary rules}

Some rules exclude highly probative evidence for reasons unrelated to accuracy: legal professional privilege protects communications between lawyer and client; the privilege against self incrimination protects a person from being compelled to give evidence against himself; the exclusionary rule in American constitutional law excludes evidence obtained by unlawful searches. These rules accept a loss of accuracy for other values: the effective provision of legal advice, the dignity of the person, the deterrence of unlawful police conduct. A decision theoretic analysis treats them as constraints that change the information set available to the fact finder. Their cost can be measured, in principle, by the expected increase in decision loss from excluding the evidence, and their benefit by their effects on behaviour. The analysis does not settle whether the trade off is worth making, but it shows that it is a trade off.

\section{Directions to the jury}

Finally, the judge's directions to the jury are a means of correcting predictable errors. English judges are guided by the Crown Court Compendium, which provides model directions on the burden and standard of proof, on identification, on bad character, on lies, on the defendant's silence, and on many other matters. The direction on lies, for example, warns the jury that people lie for many reasons, such as panic or to conceal something embarrassing, and that a lie is not evidence of guilt unless the jury is sure it was not told for an innocent reason. In Bayesian terms, the direction warns the jury that the likelihood ratio of a lie may be close to one, because innocent people lie too. The lies of Sacco and Vanzetti at their arrest, taken up in Chapter~\ref{ch:graphs}, are an example: their fear of deportation as anarchists was an alternative explanation that reduced the weight of the lies.\index{jury directions}

A Bayesian compendium would add directions on conditioning and on base rates: on the difference between the probability of the evidence and the probability of guilt, on the effect of the size of the pool of possible offenders, and on the danger of adding the weights of dependent items. Chapter~\ref{ch:program} proposes the first of these.

\section*{Notes and Questions}
\addcontentsline{toc}{section}{Notes and Questions}
\begin{enumerate}[label=\arabic*.]
\item Suppose a jury treats a likelihood ratio $\lambda$ as $\lambda^2$. For a prior probability of $0.3$ and a threshold of $0.9$, find the range of true likelihood ratios for which the distortion changes the verdict. Use your answer to frame a Rule 403 balance.
\item A defendant charged with burglary has three previous burglary convictions. Estimate the likelihood ratio of this evidence if the pool of suspects is (a) all adult males in the town, (b) men identified by searching police records for burglars, so that the record is the reason they are suspects. Discuss \emph{Hanson} in light of your answer.
\item Using Schum's three attributes of credibility alongside the four testimonial capacities, describe how cross examination reduces uncertainty about each testimonial capacity. Which hearsay exceptions address which capacities?
\item An eyewitness identifies the defendant from a fair line up of six after a ten second view at night. Assume a hit rate of $0.5$ and a false identification rate of $0.3$ spread evenly over line up members. Compute the likelihood ratio. How would it change with a suggestive line up in which the defendant alone matches the description?
\item Should Scotland abolish the corroboration requirement? Frame your answer in terms of the expected loss from false convictions and false acquittals.
\item Draft a model direction to a jury on the prosecutor's fallacy, suitable for the Crown Court Compendium.
\end{enumerate}

\section*{Further reading}
Schum's \emph{The Evidential Foundations of Probabilistic Reasoning} \citep{schum1994}; the PCAST report on forensic science \citep{pcast2016}; and, on Bayesian decision rules for the burden of persuasion, \citet{kaye1999}.
